\textbf{This section ranks nothing and recommends nothing.} It reports what the delivered files
contain and how that varies by model, by estimator and by forecast step. One quantity here is
measured against observed outcomes --- the magnitude calibration ratio, which the run itself
computed. Everything else is a property of the predictions alone. It is here because these three
behaviours determine whether a given file is being read for the thing it can support, and none of
them is visible from the files without computing it.

\subsection*{Magnitude: All Eight Under-Predict, By Factors That Differ Fourfold}

The magnitude calibration ratio is $\mathrm{MCR} = \bar{\hat{y}} / \bar{y}$, the mean prediction
over the mean observed value on the same cell-months. One is perfect aggregate calibration, below
one is systematic under-prediction. Because it averages over the draws before dividing, it is
exactly the ratio frame I's values bear to the observed totals.

\begingroup
\catcode`\_=12
\begin{table}[h!]
\centering
\small
\renewcommand{\arraystretch}{1.2}
\caption{Magnitude calibration and CRPS on the calibration test window, as computed by the run.
Ordered by state-based MCR. Lower CRPS is better; MCR nearer 1 is better calibrated.}
\label{tab:magnitude}
\begin{tabular}{@{}p{3.3cm}rrrcr@{}}
\toprule
\textbf{Model} & \textbf{MCR sb} & \textbf{MCR ns} & \textbf{MCR os} & & \textbf{CRPS sb} \\ \midrule
bright_starship  & 0.067 & 0.057 & 0.059 & & 0.1143 \\
blue_stranger    & 0.073 & 0.063 & 0.040 & & 0.1136 \\
purple_alien     & 0.093 & 0.127 & 0.095 & & 0.1130 \\
violet_visitor   & 0.110 & 0.169 & 0.106 & & 0.1132 \\
pink_pirate      & 0.195 & 0.135 & 0.060 & & 0.1115 \\
heavy_freighter  & 0.215 & 0.303 & 0.241 & & 0.1135 \\
bold_comet       & 0.240 & 0.171 & 0.257 & & 0.1107 \\
blazing_meteor   & 0.250 & 0.319 & 0.328 & & 0.1101 \\ \midrule
\textit{range}   & \textit{3.8$\times$} & \textit{5.6$\times$} & \textit{8.2$\times$} & & \textit{3.8\%} \\
\bottomrule
\end{tabular}
\end{table}
\endgroup

\noindent Two things in that table are worth stating rather than leaving to be noticed.

\textbf{Every model under-predicts magnitude, and frame I under-predicts it by between four and
fifteen times.} The least timid of the eight reaches a quarter of the observed total; the most
timid reaches a fifteenth. The shortfall is in severity rather than in location: the models mark
roughly the right places and put numbers that are too small in them. This is a property of the
estimator and the objective, not of this export --- a body trained on mean squared error over a
field that is 99.4\% zeros is rewarded for staying near zero.

\textbf{CRPS does not separate the eight the way magnitude does.} Across the eight it spans 3.8\%
while state-based magnitude spans 3.8-fold. The two orderings also run opposite ways here: the most
timid model carries the worst CRPS and the least timid the best. A 3.8\% spread is a narrow basis
for a decision, and the selection criteria --- not this note --- determine what weight it carries.

\subsection*{The Three Estimators, Per Target}

Totals over all eight models and all thirteen origins, with each frame's ratio to frame I.

\begin{table}[h!]
\centering
\small
\renewcommand{\arraystretch}{1.2}
\caption{Total predicted fatalities per frame and target, summed over 8 models $\times$ 13 origins,
with the ratio to frame I.}
\label{tab:frameratios}
\begin{tabular}{@{}lrrrrr@{}}
\toprule
 & \multicolumn{1}{c}{\textbf{I: mean}} & \multicolumn{2}{c}{\textbf{II: q95}}
 & \multicolumn{2}{c}{\textbf{III: E[y|y>0]}} \\
\cmidrule(lr){2-2}\cmidrule(lr){3-4}\cmidrule(lr){5-6}
\textbf{Target} & total & total & ratio & total & ratio \\ \midrule
State-based & 4,194,894 & 18,898,471 & 4.51$\times$ & 17,815,285 & 4.25$\times$ \\
Non-state   &   731,151 &  3,522,070 & 4.82$\times$ &  4,417,776 & 6.04$\times$ \\
One-sided   &   427,970 &  1,979,497 & 4.63$\times$ &  3,217,289 & \textbf{7.52}$\times$ \\ \midrule
All three   & 5,354,015 & 24,400,038 & 4.56$\times$ & 25,450,349 & 4.75$\times$ \\
\bottomrule
\end{tabular}
\end{table}

\noindent Frames II and III both land near five times frame I in total, which is close to the
factor by which frame I falls short of observed --- that is what makes them interesting and is the
reason all three were kept. But the per-target columns are the point of the table: frame III runs
4.25 times frame I on state-based fatalities and 7.52 times on one-sided, for the reason given in
Section~\ref{sec:frames}. A single ratio quoted for all three targets would be wrong for two of
them.

\textbf{And the two estimators are not equally well summarised by a ratio.} Taken per model and
per target --- twenty-four figures --- frame II sits between 4.25 and 4.91 times frame I, so
``about $4.6\times$'' describes every one of them. Frame III ranges from \textbf{2.88 to 11.62}
over the same twenty-four: 2.88 to 6.89 on state-based, 4.19 to 10.71 on non-state, 4.75 to 11.62
on one-sided. Its pooled 4.75 is an average over a fourfold spread and should not be used as a
conversion factor for any particular model.

\textbf{All three frames are non-zero on exactly the same rows.} 1,469,004 of the 242,678,592 rows
in each frame carry a non-zero prediction --- 0.61\%, so 99.4\% of the delivered field is exact
zeros. Those rows fall on 98,995 distinct cell-months and 6,552 distinct PRIO-GRID cells, a tenth
of the 64,818 delivered. The three estimators differ in magnitude and never in where conflict is
predicted; this was verified across all 312 files, not sampled. It is a property of the
construction rather than a coincidence: all three estimators collapse the \emph{same} draw set, and
a cell with no positive draws among its sixteen gives a mean, a 95th percentile and a conditional
mean that are all zero. Note that the figure 1,469,004 counts \emph{rows}, which means the same
cell-month is counted once per model and once per origin; it is not a count of distinct
cell-months.

\subsection*{Horizon: The Field Halves Over Five Steps, Then Stops Changing}

Total predicted state-based fatalities by forecast step, averaged over the thirteen origins and
summed over the eight models.

\begin{table}[h!]
\centering
\small
\renewcommand{\arraystretch}{1.2}
\caption{Predicted state-based fatalities by forecast step. Frame I's step-1 value is 100\%.}
\label{tab:horizon}
\begin{tabular}{@{}rrrrr@{}}
\toprule
\textbf{Step} & \textbf{I: mean} & \textbf{II: q95} & \textbf{III: E[y|y>0]}
 & \textbf{I as \% of step 1} \\ \midrule
 1 & 16,404 & 71,160 & 56,415 & 100.0\% \\
 2 & 11,883 & 52,554 & 46,616 &  72.4\% \\
 3 & 10,058 & 44,794 & 41,453 &  61.3\% \\
 4 &  9,455 & 42,590 & 39,164 &  57.6\% \\
 5 &  8,962 & 40,480 & 38,155 &  54.6\% \\
 6 &  8,789 & 39,563 & 37,213 &  53.6\% \\ \addlinespace
 9 &  8,468 & 38,259 & 36,754 &  51.6\% \\
12 &  8,393 & 37,860 & 35,606 &  51.2\% \\
18 &  8,674 & 39,266 & 37,939 &  52.9\% \\
24 &  8,607 & 38,698 & 36,794 &  52.5\% \\
30 &  8,480 & 38,353 & 36,287 &  51.7\% \\
36 &  8,728 & 39,592 & 37,336 &  53.2\% \\
\bottomrule
\end{tabular}
\end{table}

\noindent \textbf{The total settles after about five steps.} Between step 1 and step 5 the
predicted total falls by 45\%. From step 6 to step 36 --- thirty-one of the thirty-six months
delivered --- frame I stays between 8,265 and 8,907, a spread of 7.8\% with no trend, and the other
two frames behave the same way.

\textbf{What settles is the total, not the field.} These are different claims and only the first
one is true, so the distinction is worth making explicitly. Measuring how much of the field moves
between consecutive forecast steps ---
$\sum_i |f_i(h) - f_i(h-1)| \,/\, \sum_i f_i(h-1)$, the share of the predicted mass that
relocates --- gives a median of \textbf{0.93} pooled over the eight models and all thirteen
origins, and \textbf{0.94} over steps 12 to 36 alone. A frozen field would read 0. Roughly as much
of the field relocates between month 35 and month 36 as between month 1 and month 2. Per model the
figure runs from 0.67 to 1.32, and it rises slightly with the horizon rather than falling.

So the forecasts do not go flat. What goes flat is their marginal intensity: after five steps the
models predict about the same amount of conflict in total, and about the same number of active
cells, while continuing to move it around the map at full amplitude.

\textbf{What this does not establish, and it is the part that matters.} A field that keeps moving
is equally consistent with real month-to-month dynamics and with churn around a fixed
climatological intensity, and this measurement cannot separate the two. It is a property of the
predictions with no observed values in it. So nothing here licenses either reading --- neither that
month 30 carries forecast information nor that it does not. The reason to be careful with the long
horizon is Section~\ref{sec:excluded}: \textbf{nothing in this delivery has been scored at any
horizon}, and a 36-month forecast whose accuracy has not been measured past month one is untested
rather than disproven. A metric averaged evenly over the 36 months is still worth treating with
care, because 31 of those months sit in the regime where the total has stopped changing.

One further thing this does not establish. An earlier scrutiny of these eight, on a smaller region
and a shorter training budget, found per-step magnitude decaying toward zero rather than settling.
That was a different run; this measurement is of these files. The two are different quantities and
neither should be read into the other.
